% GENERATED FILE. Edit canonical lessons, then run npm run generate:books.
% canonical-source-hash: fnv1a64:ab9cc3fde8c172e4
% canonical-lessons: MR-C447-bagla, MR-C447-ghar, MR-C447-gidhad, MR-C447-kilbilne, MR-C447-sthalantar

\chapter{Herons and Kites}
\label{ch:mr-herons-and-kites}
\hlchaptermodality{447}

\begin{chapteropening}
\textbf{By the end of this chapter:} \emph{I can talk about herons, kites, vultures and migration, and the chirping of birds.}

Complete the last lesson of chapter 447: I can talk about herons, kites, vultures and migration, and the chirping of birds.
\end{chapteropening}

\section[\texorpdfstring{\mr{बगळा} (bagḷā)}{बगळा (bagḷā)}]{\mr{बगळा} (bagḷā) — a heron; an egret}

\label{lesson:MR-C447-bagla}



\begin{quote}
Before the new one: say the Marathi for to roar, then the Marathi for to graze.
\end{quote}

\subsection*{You'll want to know: \mr{बगळा}}
\textbf{\mr{बगळा}} — \emph{bagḷā} — \textquotedblleft{}a heron; an egret\textquotedblright{}.

\begin{quote}
\mr{शेतात} \mr{पांढरे} \mr{बगळे} \mr{बसले} \mr{आहेत}.
\end{quote}

\begin{quote}
\emph{śetāt pāṇḍhre bagḷe basle āhet.}
\end{quote}

\begin{quote}
White egrets are sitting in the field.
\end{quote}

Masculine. Plural: \mr{बगळे}, \emph{bagḷe}.

The plural is \mr{बगळे}.

\begin{cousinweb}
\textbf{\mr{बगळा}} (\emph{bagḷā}) belongs with Sanskrit \emph{baka}, \textquotedblleft{}a heron\textquotedblright{} (CDIAL 9115); how the Marathi form was shaped is not settled.
\end{cousinweb}

\begin{grammarlens}[title={What you've built}]
One more word for getting around.
\end{grammarlens}

\subsection*{Your turn}
\begin{itemize}
\raggedright
  \item \emph{Say it:} \emph{bagḷā}
  \item \emph{Say it:} \emph{bagḷā}, once more
  \item \emph{Say it:} \emph{śetāt pāṇḍhre bagḷe basle āhet.} — White egrets are sitting in the field.
  \item \emph{Recall:} say the Marathi for to roar, then the Marathi for to graze, then say \emph{bagḷā} again
\end{itemize}

\begin{tcolorbox}[breakable,colback=teal!4,colframe=teal!35!black,title={Before you move on}]
What does \mr{बगळा} mean? (\textquotedblleft{}A heron; an egret\textquotedblright{}.) Say it once more.
\end{tcolorbox}
\section[\texorpdfstring{\mr{घार} (ghār)}{घार (ghār)}]{\mr{घार} (ghār) — a black kite}

\label{lesson:MR-C447-ghar}



\begin{quote}
Before the new one: say the Marathi for to graze, then the Marathi for a heron.
\end{quote}

\subsection*{You'll want to know: \mr{घार}}
\textbf{\mr{घार}} — \emph{ghār} — \textquotedblleft{}a black kite\textquotedblright{}.

\begin{quote}
\mr{आकाशात} \mr{एक} \mr{घार} \mr{उडत} \mr{होती}.
\end{quote}

\begin{quote}
\emph{ākāśāt ek ghār uḍat hotī.}
\end{quote}

\begin{quote}
A kite was flying in the sky.
\end{quote}

Feminine.

A feminine noun: \mr{घार} \mr{उडत} \mr{होती}.

\begin{cousinweb}
\textbf{\mr{घार}} (\emph{ghār}) is the Marathi word for the black kite, the brown bird of prey seen over every town; its origin is not settled.
\end{cousinweb}

\begin{grammarlens}[title={What you've built}]
One more word for getting around.
\end{grammarlens}

\subsection*{Your turn}
\begin{itemize}
\raggedright
  \item \emph{Say it:} \emph{ghār}
  \item \emph{Say it:} \emph{ghār}, once more
  \item \emph{Say it:} \emph{ākāśāt ek ghār uḍat hotī.} — A kite was flying in the sky.
  \item \emph{Recall:} say the Marathi for to graze, then the Marathi for a heron, then say \emph{ghār} again
\end{itemize}

\begin{tcolorbox}[breakable,colback=teal!4,colframe=teal!35!black,title={Before you move on}]
What does \mr{घार} mean? (\textquotedblleft{}A black kite\textquotedblright{}.) Say it once more.
\end{tcolorbox}
\section[\texorpdfstring{\mr{गिधाड} (gidhāḍ)}{गिधाड (gidhāḍ)}]{\mr{गिधाड} (gidhāḍ) — a vulture}

\label{lesson:MR-C447-gidhad}



\begin{quote}
Before the new one: say the Marathi for a heron, then the Marathi for a black kite.
\end{quote}

\subsection*{You'll want to know: \mr{गिधाड}}
\textbf{\mr{गिधाड}} — \emph{gidhāḍ} — \textquotedblleft{}a vulture\textquotedblright{}.

\begin{quote}
\mr{आजकाल} \mr{गिधाडं} \mr{फार} \mr{कमी} \mr{दिसतात}.
\end{quote}

\begin{quote}
\emph{ājkāl gidhāḍaṁ phār kamī distāt.}
\end{quote}

\begin{quote}
These days vultures are seldom seen.
\end{quote}

Neuter.

A neuter noun: one \mr{गिधाड}, many \mr{गिधाडं}.

\begin{cousinweb}
\textbf{\mr{गिधाड}} (\emph{gidhāḍ}) is inherited: CDIAL lists Marathi \emph{gidhāḍ}, \textquotedblleft{}vulture, kite\textquotedblright{}, under Sanskrit \emph{gṛdhra}, \textquotedblleft{}greedy; a vulture\textquotedblright{} (CDIAL 4233).
\end{cousinweb}

\begin{grammarlens}[title={What you've built}]
One more word for getting around.
\end{grammarlens}

\subsection*{Your turn}
\begin{itemize}
\raggedright
  \item \emph{Say it:} \emph{gidhāḍ}
  \item \emph{Say it:} \emph{gidhāḍ}, once more
  \item \emph{Say it:} \emph{ājkāl gidhāḍaṁ phār kamī distāt.} — These days vultures are seldom seen.
  \item \emph{Recall:} say the Marathi for a heron, then the Marathi for a black kite, then say \emph{gidhāḍ} again
\end{itemize}

\begin{tcolorbox}[breakable,colback=teal!4,colframe=teal!35!black,title={Before you move on}]
What does \mr{गिधाड} mean? (\textquotedblleft{}A vulture\textquotedblright{}.) Say it once more.
\end{tcolorbox}
\section[\texorpdfstring{\mr{किलबिलणे} (kilbilṇe)}{किलबिलणे (kilbilṇe)}]{\mr{किलबिलणे} (kilbilṇe) — to chirp}

\label{lesson:MR-C447-kilbilne}



\begin{quote}
Before the new one: say the Marathi for a black kite, then the Marathi for a vulture.
\end{quote}

\subsection*{You'll want to know: \mr{किलबिलणे}}
\textbf{\mr{किलबिलणे}} — \emph{kilbilṇe} — \textquotedblleft{}to chirp\textquotedblright{}.

\begin{quote}
\mr{पहाटे} \mr{झाडावर} \mr{पक्षी} \mr{किलबिलतात}.
\end{quote}

\begin{quote}
\emph{pahāṭe jhāḍāvar pakṣī kilbiltāt.}
\end{quote}

\begin{quote}
At dawn the birds chirp in the tree.
\end{quote}

Past (he …): \mr{तो} \mr{किलबिलला}, \emph{to kilbillā}.

The noun is \mr{किलबिल}, the chirping of many birds.

\begin{cousinweb}
\textbf{\mr{किलबिलणे}} (\emph{kilbilṇe}) is an echo verb built on \emph{kilbil}, the twitter of small birds; its origin is not settled.
\end{cousinweb}

\begin{grammarlens}[title={What you've built}]
One more word for getting around.
\end{grammarlens}

\subsection*{Your turn}
\begin{itemize}
\raggedright
  \item \emph{Say it:} \emph{kilbilṇe}
  \item \emph{Say it:} \emph{kilbilṇe}, once more
  \item \emph{Say it:} \emph{pahāṭe jhāḍāvar pakṣī kilbiltāt.} — At dawn the birds chirp in the tree.
  \item \emph{Recall:} say the Marathi for a black kite, then the Marathi for a vulture, then say \emph{kilbilṇe} again
\end{itemize}

\begin{tcolorbox}[breakable,colback=teal!4,colframe=teal!35!black,title={Before you move on}]
What does \mr{किलबिलणे} mean? (\textquotedblleft{}To chirp\textquotedblright{}.) Say it once more.
\end{tcolorbox}
\section[\texorpdfstring{\mr{स्थलांतर} (sthalāntar)}{स्थलांतर (sthalāntar)}]{\mr{स्थलांतर} (sthalāntar) — migration}

\label{lesson:MR-C447-sthalantar}



\begin{quote}
Before the new one: say the Marathi for a vulture, then the Marathi for to chirp.
\end{quote}

\subsection*{You'll want to know: \mr{स्थलांतर}}
\textbf{\mr{स्थलांतर}} — \emph{sthalāntar} — \textquotedblleft{}migration\textquotedblright{}.

\begin{quote}
\mr{हिवाळ्यात} \mr{अनेक} \mr{पक्षी} \mr{इथे} \mr{स्थलांतर} \mr{करतात}.
\end{quote}

\begin{quote}
\emph{hivāḷyāt anek pakṣī ithe sthalāntar kartāt.}
\end{quote}

\begin{quote}
In winter many birds migrate here.
\end{quote}

Neuter.

\mr{स्थलांतर} \mr{करणे}, to migrate; also people moving for work.

\begin{cousinweb}
\textbf{\mr{स्थलांतर}} (\emph{sthalāntar}) is a tatsama, taken from Sanskrit \emph{sthalāntara}, \textquotedblleft{}another place\textquotedblright{} (Monier-Williams): going over to another place.
\end{cousinweb}

\begin{grammarlens}[title={What you've built}]
One more word for getting around.
\end{grammarlens}

\subsection*{Your turn}
\begin{itemize}
\raggedright
  \item \emph{Say it:} \emph{sthalāntar}
  \item \emph{Say it:} \emph{sthalāntar}, once more
  \item \emph{Say it:} \emph{hivāḷyāt anek pakṣī ithe sthalāntar kartāt.} — In winter many birds migrate here.
  \item \emph{Recall:} say the Marathi for a vulture, then the Marathi for to chirp, then say \emph{sthalāntar} again
\end{itemize}

\begin{tcolorbox}[breakable,colback=teal!4,colframe=teal!35!black,title={Before you move on}]
What does \mr{स्थलांतर} mean? (\textquotedblleft{}Migration\textquotedblright{}.) Say it once more.
\end{tcolorbox}
